\chapter*{Conclusion Générale}
\label{chapter:Conclusion}
\addcontentsline{toc}{chapter}{Conclusion Générale}
\markboth{Conclusion Générale}{}
\setcounter{section}{0}

Ce rapport a présenté les travaux réalisés durant le stage de perfectionnement effectué
au sein du département Informatique (IT) du Mövenpick Hôtel de Tanger, portant sur la
conception et le développement de \textbf{NetSentry}, un outil interne de découverte et
de cartographie des équipements connectés au réseau local de l'hôtel.

\section{Bilan des réalisations}

Le premier chapitre a permis de poser le contexte général du stage et du projet, en
mettant en évidence l'absence, au sein du département IT, d'un inventaire fiable des
équipements connectés au réseau, et a montré, à travers une étude de l'existant et une
analyse SWOT, qu'aucune solution disponible sur le marché ne combinait simplicité de
déploiement, auto-hébergement et historique exploitable pour ce périmètre restreint,
justifiant ainsi le développement d'une solution sur mesure. Le deuxième chapitre a
détaillé les exigences retenues puis formalisé la conception du système à travers ses
diagrammes UML. Le troisième chapitre a détaillé l'architecture technique retenue --
avec Nmap comme moteur de scan -- ainsi que les technologies mobilisées. Le quatrième
chapitre a enfin illustré, par des captures d'écran, le fonctionnement réel et vérifié de
l'outil : découverte des équipements, tableau de bord, alertes, historique des scans et
déploiement conteneurisé.

\section{Défis rencontrés et solutions apportées}

Plusieurs défis ont marqué la réalisation de ce projet :
\begin{itemize}
    \item \textbf{Périmètre accessible limité} : les systèmes métiers de l'hôtel (PMS,
    point de vente) étant gérés par un prestataire externe, le sujet initial a dû être
    recentré sur le seul réseau interne, réellement accessible au département IT.
    \item \textbf{Fiabilité de l'identification des équipements} : certains appareils
    (notamment les smartphones modernes) utilisent des adresses MAC aléatoires à des
    fins de confidentialité, rendant leur fabricant impossible à déterminer -- une
    limitation inhérente à toute identification basée sur l'adresse MAC, documentée plutôt
    que masquée dans l'outil.
    \item \textbf{Contrainte de simplicité} : concevoir un tableau de bord exploitable
    par du personnel non spécialisé en réseau, tout en s'appuyant sur un outil aussi
    riche que Nmap, a nécessité de limiter volontairement certaines fonctionnalités
    avancées (détection de système d'exploitation notamment).
\end{itemize}

\section{Compétences acquises et renforcées}

Ce stage a permis de renforcer plusieurs compétences clés :
\begin{itemize}
    \item \textbf{Compétences réseau} : adressage IP/MAC, protocole ARP, notion de
    ports, et prise en main approfondie de Nmap.
    \item \textbf{Compétences techniques} : développement backend (Node.js, Express,
    PostgreSQL), frontend (React), conteneurisation (Docker), et sécurisation d'une API
    (authentification JWT).
    \item \textbf{Modélisation et architecture} : conception UML (cas d'utilisation,
    séquence, modèle conceptuel de données) et mise en \oe{}uvre d'une architecture en
    couches.
    \item \textbf{Compétences transversales} : compréhension du fonctionnement d'un
    département informatique en entreprise et de sa collaboration avec un prestataire
    externe, autonomie et adaptation du périmètre d'un projet aux contraintes réelles
    d'une organisation.
\end{itemize}

\section{Perspectives d'évolution}

Le système développé constitue une base solide pour plusieurs évolutions futures :

\begin{itemize}
    \item \textbf{Notifications par e-mail} : informer automatiquement l'administrateur
    IT dès la détection d'un équipement inconnu, en complément de l'alerte affichée dans
    le tableau de bord.
    \item \textbf{Classification par type d'appareil} : enrichir l'identification des
    équipements par une heuristique combinant fabricant et ports ouverts.
    \item \textbf{Historique Git plus granulaire} : poursuivre le versionnage du projet
    par des commits plus fréquents au fil des prochaines itérations.
    \item \textbf{Déploiement en production} : installation de NetSentry sur le serveur
    définitif de l'hôtel, dans les conditions réelles d'exploitation du réseau.
\end{itemize}

En conclusion, ce stage a représenté une étape importante dans le parcours de
l'étudiant, en le confrontant aux réalités d'un développement logiciel mené en
autonomie, à la croisée du développement web et de l'administration réseau, au service
d'un besoin réel exprimé par une entreprise d'accueil.
